\documentclass[11pt]{article}
\usepackage[margin=1in]{geometry}
\usepackage{amsmath}
\usepackage{booktabs}
\usepackage{url}
\title{Synchronizing the manuscript with ea88059}
\date{Prepared October 11, 2026}
\begin{document}
\maketitle

\section{Question}
The manuscript \path{paper/manuscript.tex} described the controller of
October 5 (\texttt{d1c31d6}): a separating, road-recoverable terminal set
that was not invariant, a CLF certified only for the linearized model with
its gain discarded, a solution issued directly, and 82 of 84 noisy
encounters recovered. Since then the controller changed in eleven commits
(\texttt{68883d2} to \texttt{ea88059}): the issued plan is the nonlinear
rollout of an accepted step; the terminal set is the braking continuation of
both vehicles, holding only mild endpoints, with the obstacle's braking
ending at rest; the vehicle model has a kinematic regime below 4~m/s so that
it can stop; the cruise law $u=u^*+Ke$ is certified with the CLF on the
sampled nonlinear model; and the problem is assembled about twice as fast
with looser solver tolerances. The user asked for the paper to be brought in
line with the algorithm.

\section{What was read}
The manuscript line by line; the commit history since \texttt{4ccf41f};
\path{controller/PCBF_CLF_ARCHITECTURE.md}, \path{controller/TERMINAL_SAFE_SET.md},
\path{controller/NOMINAL_CLF.md}, \path{controller/CONTROLLER_FILES.md};
\path{controller/solvePredictiveControl.m}, \path{controller/terminalSafeSet.m},
\path{controller/nonlinearBicycleModel.m}, \path{controller/collisionAvoidanceController.m},
the target forecast of \path{controller/predictiveSafetyGeometry.m},
\path{config/collisionAvoidanceControllerConfig.m}, \path{config/clfMatrices.json}
and the header of \path{scripts/synthesizeClfMatrices.m}; the reports of
October 11 (\path{BRAKING_TERMINAL_SET_20261011.tex},
\path{TARGET_STOP_AT_REST_20261011.tex}, \path{MILD_TERMINAL_SET_20261011.tex},
\path{LOW_SPEED_MODEL_20261011.tex}, \path{CONTROLLER_SPEED_20261011.tex}) and the
index; and the estimator document changes since \texttt{d1c31d6}, which are
consolidations and notes without algorithm changes. Statements taken from
the documents were checked against the code where they enter the manuscript.

\section{Manuscript changes}
\begin{itemize}
\item Abstract, contributions and the boundary paragraph of the
  introduction, the architecture paragraph of Section II and the
  safety-argument subsection: the terminal set from the obstacle's model,
  the issued plan as a rollout, recursive feasibility proved on the declared
  model and measured with the estimator, the cruise law certified on the
  sampled nonlinear model, 74 of 84.
\item Section II: the tire model is used above 4~m/s; braking of the
  obstacle ends at rest. Section IV-C: the force-balance premise is the Fiala
  model, which the plant follows exactly above 4~m/s only.
\item Section V, rewritten: the low-speed regime of the prediction model
  (kinematic bicycle, blend, exact rest) and the obstacle forecast with rest;
  the anchor (startup and onset rollouts, the shifted issued plan and its
  extension by cruise, guidance and braking holds); collision rows
  unchanged; the handling rows' speed limits; the terminal set with the
  terminal law, mildness, Assumption~3, the closed-form continuation, the
  parked check and Proposition~1 (invariance); the CLF with the ray-average
  identity, the zonotope enclosure, the LMI with Petersen's lemma, the
  self-consistent growth and the enclosure hypothesis; the enclosure use with
  the terminal widening; the two programs with the tolerances, the trust
  expansion, the fresh rollout and the step-size rule; Proposition~2
  (recursive feasibility under H1--H3); the trust scale from the full step's
  remainder; the scope; and the computation measures that leave the problem
  unchanged. The separating proposition, the road-recovery cone and the
  innovation split are removed.
\item Section VI: the parameter table (kinematic regime, $a_b$, $a_m$,
  certification box and input bound, solver tolerances; the retired
  $v_{\rm sep}$, $c_{\rm rec}$, $R_V$, $c_V$, $v_f$ removed); the certified sets
  and contractions; the plant's low-speed regime; the outcomes table and the
  figure for the new campaign; a new subsection on the shifted plan and the
  terminal set in the loop; the estimation audit of the new campaign; the ten
  failures with a table of the published obstacle parameters; computation.
\item Discussion and conclusion rewritten; the AI-assistance disclosure
  extended.
\item Four references added: Chen et al.\ (2021, backup CBFs), Petersen
  (1987), Chisci et al.\ (2001), Mayne et al.\ (2005).
\end{itemize}

\section{Sources of the numbers}
No controller simulation and no unit test was run for this revision.
\begin{itemize}
\item The campaign is the reference arm of a solver-tolerance study that the
  user was running in another session on October 11 (raw outputs under
  \path{~/.cache/collisionAvoidance/solver_tolerance_20261011/out-ref}, a
  \texttt{git archive} of \texttt{ea88059} with a solver link, verified
  identical to the commit's \path{config}, \path{controller}, \path{estimator}
  and \path{scripts}). Its 98 runs reproduce the committed campaign of
  \path{CONTROLLER_SPEED_20261011/campaign-runs.csv} in outcome, frame count,
  least clearance, largest lateral deviation and recovery time, run for run.
  Six MATLAB sessions shared the machine; the timings quoted are those of
  this rerun, so that the figure and the statistics come from one record.
\item Outcomes, gaps, road margins, recovery times, extremes, horizon
  statistics, accepted steps, shifted-candidate feasibility, extension
  counts, continuation ends, primary optima, CLF decrease counts, estimation
  errors and target-parameter errors: derived from the traces by
  \path{MANUSCRIPT_ALGORITHM_SYNC_20261011/deriveClosedLoopStatistics.py}
  (output \path{closed-loop-statistics.json}). The CLF counts restrict the
  pairs to $V>10^{-3}$ and to samples whose full step was accepted. The
  cruise statistics of the exact 15~m/s curved head-on (speed error, lateral
  offset, slack frames between 20 and 50~s) and the forecast displacements of
  the failure table were computed from the same traces with session scripts
  that reproduce the target forecast of \texttt{predictTarget}.
\item Estimator audit: \path{scripts/replayEstimatorValidity.m} of
  \texttt{ea88059} on the seed-20261003 runs, through
  \path{MANUSCRIPT_ALGORITHM_SYNC_20261011/auditEstimatorInLoop.m}; output
  \path{estimator-audit-seed20261003.json} and \path{.txt}. 4001 frames, ego
  bound valid and containing the truth in all, target bound available in
  2018 and containing the truth in all, no frame outside the sideslip cone or
  the adhesion condition, replay difference at most $5\cdot10^{-15}$.
\item Certified sets, contractions, hold factors and the sampled check:
  \path{config/clfMatrices.json} (entries of October 10, certificate seed
  20261010).
\item The benchmark with a target that appears, the previous
  implementation's timings, the tolerance study and the offline re-solves of
  four failing samples: \path{CONTROLLER_SPEED_20261011.tex}.
\item The kinematic comparison of the October 5 estimator (heading errors
  0.07--0.31~rad, five of fourteen recovered): \path{JOINT_DESIGN_20261005.tex},
  quoted as a result on the controller of that date.
\item The full-horizon target tube measurement (31 of 84, one collision):
  \path{TARGET_PARAMETER_SET_20261008.tex}; the set-membership enclosure of
  the obstacle's constants mentioned in the discussion: the same report.
\item Figure \path{closed_loop_results_v03}:
  \path{MANUSCRIPT_ALGORITHM_SYNC_20261011/figure-sources/closed_loop_results_v03.py}.
\end{itemize}

\section{Documentation changes}
\path{paper/README.md} and \path{report/README.md} (this entry). No
executable statement and no controller document was changed.

\section{Open items for review}
\begin{itemize}
\item Section III-B (road-curb extraction) describes the offline
  curb-detection experiment that was removed from the repository on
  October 6, 2026 at the user's direction; its figure cannot be regenerated
  from the repository. Whether the subsection stays is the user's decision;
  it was left unchanged here.
\item A concurrent session committed, while this revision was being
  written, a floor on the tie term's weight on the shifted plan with the
  solver-tolerance study (\texttt{e207dde},
  \path{SOLVER_TOLERANCE_20261011.tex}: cruise exact at both speeds, 14 of 14
  and 73 of 84). Its commit also carried this revision's files, which were
  staged in the shared index at the time, so the manuscript revision is
  recorded under that commit's message; the follow-up commit that records
  it is the one this report is cited from. The manuscript describes
  \texttt{ea88059} and its campaign, and its discussion mentions the
  later floor; the explanation it offers for the CLF increases at 15~m/s
  on the curve (the tie term's weight vanishes at the path) follows from
  the design as written and agrees with that session's finding.
\item The derivations of Section V restate the repository documents and
  code; the proofs of the two propositions and the Petersen form of the LMI
  have not been independently reviewed. Assumption~3 is the user's modeling
  assumption and carries no mismatch bound.
\item The bibliographic data of the four added references should be checked
  against the publishers' records.
\item The estimator-alone comparison (Sections VI-A and VI-B) was not rerun.
\item The timings are those of six runs sharing the machine; the
  lighter-load benchmark quoted alongside is from the speed report.
\item The architecture drawing still labels the road-boundary branch ``not
  yet used''; the caption says so.
\item Section VI-F quotes the offline re-solves of four failing samples from
  the speed report; the remaining six failing samples were not re-solved.
\end{itemize}
\end{document}
